% nats-relay-passA-nowedge.tex
% FROZEN PASS-A ARTIFACT --- DO NOT EDIT except in lockstep with
% docs/nats-relay.tex.  This is docs/nats-relay.tex with the single
% WedgeDetected recovery operation removed, modelling the PRE-FIX code
% (nats-py keepalive default of 240s that never fired in practice).
% Running `z-spec test` on this file MUST fail with a reachable
% terminal halfOpen state: that deadlock IS biff-tww, proven from the
% model.  See DESIGN.md DES-041.
%
\documentclass[a4paper,10pt,fleqn]{article}
\usepackage[margin=1in]{geometry}
\usepackage{fuzz}
\usepackage[colorlinks=true,linkcolor=blue,citecolor=blue,urlcolor=blue]{hyperref}

\begin{document}

\title{Biff NATS Relay: A Z Specification}
\author{Formal Model of Connection Lifecycle and KV Watcher}
\date{March 2026}
\hypersetup{
  pdftitle={Biff NATS Relay: A Z Specification},
  pdfauthor={punt-labs},
  pdfsubject={Z Specification},
  pdfcreator={fuzz/probcli}
}
\maketitle

\tableofcontents
\newpage

%%----------------------------------------------------------------------
\section{Introduction}
%%----------------------------------------------------------------------

The NATS relay is the networked backend for biff's communication
system.  It manages a TCP connection to a NATS server, lazily
provisions infrastructure (KV bucket, JetStream streams), and runs
a KV watcher that detects session changes and wall broadcasts in
real-time.

This specification models three interacting state machines:

\begin{enumerate}
  \item \textbf{Connection lifecycle} --- lazy connect, disconnect,
    reconnect, and close transitions.
  \item \textbf{KV watcher} --- the long-lived watcher that observes
    session and wall changes, survives snapshot-done markers and
    timeouts, and restarts only on errors.
  \item \textbf{Session presence} --- KV-backed sessions with TTL
    expiry, heartbeat refresh, and logout detection.
\end{enumerate}

The key invariant proven by this model: the KV watcher remains
alive across snapshot-done markers and idle timeouts.  Only
explicit shutdown or unrecoverable errors terminate it.  This
captures the fix for biff-udp.

A second invariant, added to close a recurring class of wedge
bugs, concerns the connection itself.  A NATS connection can
enter a \emph{half-open} state: the TCP socket is up
(\texttt{is\_closed} is false) but every JetStream or KV request
times out.  If no transition can leave that state, the relay is
wedged --- every tool call blocks on \texttt{nats: timeout} with
no recovery (biff-tww).  This specification adds the half-open
state explicitly and proves a liveness property: the connection
can never be stuck there.  The proof is by deadlock analysis ---
in the absence of a wedge-detection transition, probcli finds
half-open to be a reachable terminal state, and that terminal
state \emph{is} biff-tww.  With the recovery transition present,
half-open is deadlock-free.

%%----------------------------------------------------------------------
\section{Basic Types}
%%----------------------------------------------------------------------

\begin{zed}
  [USER, TTY, REPO, KVKEY, KVVALUE, SESSIONKEY]
\end{zed}

$USER$ identifies a human user.  $TTY$ identifies a terminal
session (random 8-char hex).  $REPO$ scopes all state to a
repository.  $KVKEY$ and $KVVALUE$ are opaque KV store keys and
values.  $SESSIONKEY$ is the composite $\{user\}:\{tty\}$ identifier.

%%----------------------------------------------------------------------
\section{Free Types}
%%----------------------------------------------------------------------

\begin{zed}
  ZBOOL ::= ztrue | zfalse
\end{zed}

\begin{zed}
  ConnState ::= disconnected | reconnecting | connected | halfOpen | closed
\end{zed}

The connection has five states.  $disconnected$ is a deliberate,
reversible release of the TCP connection (the relay's own
\texttt{disconnect()}); a subsequent call reconnects.
$reconnecting$ is distinct: nats-py's background reconnect loop is
in progress after a dropped socket --- \texttt{is\_closed} is
false, \texttt{is\_connected} is false, and cached handles have
already been invalidated by \texttt{\_on\_disconnect}.
$connected$ is an active TCP connection whose JetStream and KV
handles are provisioned.  $halfOpen$ is the wedge: the socket is
up (\texttt{is\_closed} false) but every JetStream or KV request
times out; the cached handles are stale and must not be used.
$closed$ is permanent shutdown --- no reconnection is possible.

\begin{zed}
  ReqOutcome ::= reqOk | reqTimeout
\end{zed}

The outcome of a JetStream or KV request issued on a $connected$
connection.  $reqOk$ is a normal response.  $reqTimeout$ is a
\texttt{nats.errors.TimeoutError}: the request was sent but no
reply arrived within the deadline, which is the observable
signature of a half-open connection.

\begin{zed}
  WatcherPhase ::= idle | snapshotting | streaming | stopped
\end{zed}

The KV watcher progresses through phases: $idle$ (not yet started),
$snapshotting$ (receiving the initial KV snapshot), $streaming$
(receiving live updates post-snapshot), and $stopped$ (shutdown
complete).

\begin{zed}
  KVOp ::= kvPut | kvDelete | kvPurge
\end{zed}

Operations observed on KV entries.  $kvPut$ stores or updates a
value.  $kvDelete$ and $kvPurge$ remove entries (TTL expiry
produces $kvDelete$).

\begin{zed}
  WatcherEvent ::=
  weEntry | weSnapshotDone | weTimeout | weError | weShutdown
\end{zed}

Events that the KV watcher loop must handle.  $weEntry$ is a KV
change.  $weSnapshotDone$ is the $None$ marker from
\texttt{watcher.updates()}.  $weTimeout$ is a
\texttt{TimeoutError} from idle periods.  $weError$ is an
unrecoverable NATS error.  $weShutdown$ is the shutdown signal.

%%----------------------------------------------------------------------
\section{Global Constants}
%%----------------------------------------------------------------------

\begin{axdef}
  kvTtlSeconds : \nat \\
  watcherTimeoutSeconds : \nat \\
  maxSessions : \nat \\
  restartDelaySeconds : \nat
  \where
  kvTtlSeconds = 259200 \\
  watcherTimeoutSeconds = 5 \\
  maxSessions = 1000 \\
  restartDelaySeconds = 2
\end{axdef}

$kvTtlSeconds$ is the NATS KV TTL (3 days).
$watcherTimeoutSeconds$ is the \texttt{updates()} timeout that
keeps the shutdown check responsive.  $maxSessions$ bounds the
model for animation.

%%----------------------------------------------------------------------
\section{Connection State}
%%----------------------------------------------------------------------

\begin{schema}{Connection}
  connState : ConnState \\
  kvProvisioned : ZBOOL \\
  streamsProvisioned : ZBOOL \\
  namesProvisioned : ZBOOL \\
  wtmpAvailable : ZBOOL
  \where
  connState \neq connected \implies
  kvProvisioned = zfalse \land
  streamsProvisioned = zfalse \land
  namesProvisioned = zfalse \\
  connState = connected \implies
  kvProvisioned = ztrue \land
  streamsProvisioned = ztrue \land
  namesProvisioned = ztrue
\end{schema}

This is the generalised handle-validity invariant: valid
JetStream, sessions-KV, and names-KV handles exist \emph{if and
only if} the connection is $connected$.  In particular they do
not exist in $halfOpen$ --- the stale handles cached during the
wedge are, in the model, no handles at all.  This generalises the
earlier ``$disconnected$ or $closed \implies$ no handles'' (DES-029)
to cover $reconnecting$ and $halfOpen$ as well.  It is the
formal statement that a half-open connection must never be treated
as usable: the biff-tww crash-loop arose precisely because the code
kept and re-used handles that this invariant forbids.
$wtmpAvailable$ degrades independently --- wtmp provisioning
failure is non-fatal --- and is cleared whenever the handles are.

\begin{schema}{InitConnection}
  Connection'
  \where
  connState' = disconnected \\
  kvProvisioned' = zfalse \\
  streamsProvisioned' = zfalse \\
  namesProvisioned' = zfalse \\
  wtmpAvailable' = zfalse
\end{schema}

%%----------------------------------------------------------------------
\section{Connection Operations}
%%----------------------------------------------------------------------

\subsection{Ensure Connected (Lazy Connect)}

\begin{schema}{EnsureConnected}
  \Delta Connection
  \where
  connState = disconnected \lor connState = connected \\
  connState' = connected \\
  kvProvisioned' = ztrue \\
  streamsProvisioned' = ztrue \\
  namesProvisioned' = ztrue \\
  wtmpAvailable' = wtmpAvailable
\end{schema}

The relay's lazy connect (\texttt{\_ensure\_connected}).  It is
enabled only from $disconnected$ (a fresh connect) or $connected$
(idempotent no-op).  It is \emph{not} enabled from $halfOpen$:
this is the formal counterpart of the DES-029 bug.  The fast path
sees \texttt{is\_closed} false and returns the stale cached
handles without reconnecting, so $EnsureConnected$ makes no
progress out of the wedge and must not be modelled as if it did.
Nor is it enabled from $reconnecting$ (nats-py owns that window;
provisioning would block until reconnect completes, bounded by
\texttt{\_CONNECT\_PROVISION\_TIMEOUT}, biff-wr3) or from $closed$
(no reconnection after close).  $wtmpAvailable$ is preserved ---
it only changes during provisioning.

\subsection{Disconnect (Reversible)}

\begin{schema}{Disconnect}
  \Delta Connection
  \where
  connState = connected \\
  connState' = disconnected \\
  kvProvisioned' = zfalse \\
  streamsProvisioned' = zfalse \\
  namesProvisioned' = zfalse \\
  wtmpAvailable' = zfalse
\end{schema}

The relay's explicit, reversible \texttt{disconnect()}: it acts
on a healthy $connected$ connection and releases it.  The next
relay call reconnects via $EnsureConnected$.  It is deliberately
\emph{not} the escape from $halfOpen$ --- the running server does
not autonomously call \texttt{disconnect()} while wedged; only
keepalive-driven wedge detection ($WedgeDetected$) can leave that
state.

\subsection{Close (Permanent)}

\begin{schema}{Close}
  \Delta Connection
  \where
  connState \neq halfOpen \\
  connState' = closed \\
  kvProvisioned' = zfalse \\
  streamsProvisioned' = zfalse \\
  namesProvisioned' = zfalse \\
  wtmpAvailable' = zfalse
\end{schema}

Permanent shutdown, reachable from any \emph{definite} transport
state.  The precondition $connState \neq halfOpen$ records a
modelling decision that is load-bearing for the liveness proof:
an orderly close cannot be the escape from the wedge.  In the
biff-tww crash-loop the server never reaches an orderly
\texttt{close()} --- it is blocked awaiting a request that will
never return --- so $halfOpen$ must be escaped by autonomous
wedge detection, not by $Close$.  From $closed$ the operation is
an idempotent self-loop, which keeps $closed$ deadlock-free.

\subsection{Reset Infrastructure}

\begin{schema}{ResetInfrastructure}
  \Delta Connection
  \where
  connState \neq halfOpen \\
  connState' = connState \\
  kvProvisioned' = zfalse \\
  streamsProvisioned' = zfalse \\
  namesProvisioned' = zfalse \\
  wtmpAvailable' = zfalse
\end{schema}

Clears cached handles without changing the transport state
(\texttt{reset\_infrastructure()}).  By the handle-validity
invariant, $connState' = connState$ together with cleared handles
forces $connState \neq connected$; the operation is therefore a
self-loop at $disconnected$, $reconnecting$, and $closed$, where
the handles are already absent.  It is excluded from $halfOpen$
for the same reason as $Close$: clearing a cached handle does not
lift the wedge (the socket is still up and still unresponsive),
so it must not be modelled as an escape.

\subsection{Request (JetStream / KV Round-Trip)}

\begin{schema}{Request}
  \Delta Connection \\
  result! : ReqOutcome
  \where
  connState = connected \\
  ( result! = reqOk \land
    connState' = connected \land
    kvProvisioned' = ztrue \land
    streamsProvisioned' = ztrue \land
    namesProvisioned' = ztrue \land
    wtmpAvailable' = wtmpAvailable ) \\
  \lor \\
  ( result! = reqTimeout \land
    connState' = halfOpen \land
    kvProvisioned' = zfalse \land
    streamsProvisioned' = zfalse \land
    namesProvisioned' = zfalse \land
    wtmpAvailable' = zfalse )
\end{schema}

Every JetStream or KV round-trip issued by the relay ---
\texttt{stream\_info} in \texttt{get\_unread\_summary}, \texttt{kv.get}
in \texttt{heartbeat}, and the like.  From $connected$ it has two
outcomes.  $reqOk$ leaves the connection $connected$ with its
handles intact.  $reqTimeout$ is the half-open signature: the
request timed out, the connection moves to $halfOpen$, and the
handles are invalidated to satisfy the handle-validity invariant.
This is the operation that makes $halfOpen$ reachable, and the
crash-loop of biff-tww is exactly $reqTimeout$ fired repeatedly
with no way out.

\subsection{Reconnect (nats-py Auto-Reconnect Completes)}

\begin{schema}{Reconnect}
  \Delta Connection
  \where
  connState = reconnecting \\
  connState' = connected \\
  kvProvisioned' = ztrue \\
  streamsProvisioned' = ztrue \\
  namesProvisioned' = ztrue \\
  wtmpAvailable' = wtmpAvailable
\end{schema}

nats-py's background reconnect succeeds and fires
\texttt{\_on\_reconnect}; the next relay call re-provisions the
handles.  This closes the recovery path
$connected \to halfOpen \to reconnecting \to connected$, so no
reachable state is stuck in $halfOpen$.

%%----------------------------------------------------------------------
\section{Session Presence}
%%----------------------------------------------------------------------

\begin{schema}{SessionStore}
  sessions : SESSIONKEY \pinj KVVALUE \\
  wallValue : KVKEY \pfun KVVALUE
  \where
  \# sessions \leq maxSessions \\
  \# wallValue \leq 1
\end{schema}

Sessions are a partial injection: each session key maps to
exactly one value (the serialised $UserSession$).  $wallValue$
is either empty (no wall) or a single mapping (active wall).
Initially the store is empty: $sessions = \emptyset$ and
$wallValue = \emptyset$.  There is no separate $InitSessionStore$
schema: the model-checked state is the $Connection$ machine (see
\S\ref{sec:modelcheck}), and the session store is exercised by
the operations below only while $connState = connected$.

\subsection{Update Session (Heartbeat / Login)}

\begin{schema}{UpdateSession}
  \Delta SessionStore \\
  \Xi Connection \\
  key? : SESSIONKEY \\
  value? : KVVALUE
  \where
  connState = connected \\
  sessions' = sessions \oplus \{ key? \mapsto value? \} \\
  wallValue' = wallValue \\
  \# sessions' \leq maxSessions
\end{schema}

Creates or refreshes a session.  Override ($\oplus$) ensures
idempotent updates.

\subsection{Delete Session (Logout / TTL Expiry)}

\begin{schema}{DeleteSession}
  \Delta SessionStore \\
  \Xi Connection \\
  key? : SESSIONKEY
  \where
  connState = connected \\
  key? \in \dom sessions \\
  sessions' = \{ key? \} \ndres sessions \\
  wallValue' = wallValue
\end{schema}

Removes a session.  In production, this is triggered by
explicit delete or KV TTL expiry.

\subsection{Set Wall}

\begin{schema}{SetWall}
  \Delta SessionStore \\
  \Xi Connection \\
  wallKey? : KVKEY \\
  wallVal? : KVVALUE
  \where
  connState = connected \\
  sessions' = sessions \\
  wallValue' = \{ wallKey? \mapsto wallVal? \}
\end{schema}

\subsection{Clear Wall}

\begin{schema}{ClearWall}
  \Delta SessionStore \\
  \Xi Connection
  \where
  connState = connected \\
  sessions' = sessions \\
  wallValue' = \emptyset
\end{schema}

%%----------------------------------------------------------------------
\section{KV Watcher State Machine}
%%----------------------------------------------------------------------

This section models the watcher loop from \texttt{\_run\_kv\_watch}
and its supervisor \texttt{\_kv\_watcher\_loop}.

\begin{schema}{Watcher}
  phase : WatcherPhase \\
  shutdownRequested : ZBOOL \\
  restartCount : \nat
  \where
  restartCount \leq 100 \\
  phase = stopped \implies shutdownRequested = ztrue
\end{schema}

The watcher tracks its current phase, whether shutdown has been
requested, and how many times the supervisor has restarted it.
A stopped watcher implies shutdown was requested.  Initially the
watcher is $idle$, with $shutdownRequested = zfalse$ and
$restartCount = 0$.  As with the session store, there is no
separate $InitWatcher$ schema; the watcher operations are enabled
only while $connState = connected$, which is consistent with the
composed invariant $phase \neq idle \implies connState = connected$.

\subsection{Start Watching}

\begin{schema}{StartWatching}
  \Delta Watcher \\
  \Xi Connection
  \where
  connState = connected \\
  phase = idle \\
  shutdownRequested = zfalse \\
  phase' = snapshotting \\
  shutdownRequested' = zfalse \\
  restartCount' = restartCount
\end{schema}

Transitions from idle to snapshotting when the watcher is
created.  The first entries from \texttt{watchall()} are the
snapshot.

\subsection{Snapshot Done}

\begin{schema}{SnapshotDone}
  \Delta Watcher \\
  \Xi Connection
  \where
  connState = connected \\
  phase = snapshotting \\
  shutdownRequested = zfalse \\
  phase' = streaming \\
  shutdownRequested' = zfalse \\
  restartCount' = restartCount
\end{schema}

The $None$ marker from \texttt{updates()} signals the end of
the snapshot.  \textbf{The watcher continues into streaming
phase.}  This is the key property: the old \texttt{async~for}
pattern terminated here (biff-udp).

\subsection{Receive Entry}

\begin{schema}{ReceiveEntry}
  \Delta Watcher \\
  \Xi Connection
  \where
  connState = connected \\
  (phase = snapshotting \lor phase = streaming) \\
  shutdownRequested = zfalse \\
  phase' = phase \\
  shutdownRequested' = zfalse \\
  restartCount' = restartCount
\end{schema}

A KV entry is received.  The phase does not change --- entries
arrive during both snapshotting and streaming.

\subsection{Handle Timeout}

\begin{schema}{HandleTimeout}
  \Delta Watcher \\
  \Xi Connection
  \where
  connState = connected \\
  phase = streaming \\
  shutdownRequested = zfalse \\
  phase' = streaming \\
  shutdownRequested' = zfalse \\
  restartCount' = restartCount
\end{schema}

A \texttt{TimeoutError} from \texttt{updates(timeout=5.0)} is
caught and the loop continues.  \textbf{The watcher stays in
streaming phase.}  Without the catch, this would propagate to
the supervisor and trigger a restart.

\subsection{Handle Error (Supervisor Restart)}

\begin{schema}{HandleError}
  \Delta Watcher \\
  \Xi Connection
  \where
  connState = connected \\
  phase \neq idle \\
  phase \neq stopped \\
  shutdownRequested = zfalse \\
  phase' = snapshotting \\
  shutdownRequested' = zfalse \\
  restartCount' = restartCount + 1
\end{schema}

An unrecoverable error causes the supervisor to restart the
watcher.  The new watcher begins in snapshotting phase
(a fresh \texttt{watchall()}).

\subsection{Request Shutdown}

\begin{schema}{RequestShutdown}
  \Delta Watcher \\
  \Xi Connection
  \where
  connState = connected \\
  shutdownRequested = zfalse \\
  shutdownRequested' = ztrue \\
  phase' = stopped \\
  restartCount' = restartCount
\end{schema}

The shutdown event signals the watcher to exit.  The loop
checks \texttt{shutdown.is\_set()} on every iteration.

%%----------------------------------------------------------------------
\section{Composed System State}
%%----------------------------------------------------------------------

\begin{schema}{NatsRelayState}
  Connection \\
  SessionStore \\
  Watcher
  \where
  phase \neq idle \implies connState = connected \\
  phase = streaming \implies kvProvisioned = ztrue
\end{schema}

The watcher can only be active when the connection is
established.  Streaming implies KV is provisioned (the watcher
reads from it).  With the generalised handle-validity invariant,
$phase = streaming \implies kvProvisioned = ztrue$ also implies
$phase = streaming \implies connState = connected$: a wedge
($halfOpen$, where $kvProvisioned = zfalse$) cannot coexist with a
streaming watcher, which is consistent --- a wedge manifests to
the watcher as a $weError$ and drives it through $HandleError$.

The composed initial state is: $connState = disconnected$ with no
handles ($InitConnection$), an empty session store, and an $idle$
watcher.  $InitConnection$ is the sole initialisation schema so
that the model checker explores the $Connection$ lifecycle as a
single, well-posed state machine; the session and watcher
components are exercised as connection-gated operations rather
than as independent initial states (see \S\ref{sec:modelcheck}).

%%----------------------------------------------------------------------
\section{Key Properties and Model Checking}
\label{sec:modelcheck}
%%----------------------------------------------------------------------

The specification captures these properties, each of which can be
verified by model checking with probcli (\texttt{z-spec test}):

\begin{enumerate}
  \item \textbf{Handle validity}: valid JetStream, sessions-KV, and
    names-KV handles exist if and only if $connState = connected$.
    A $halfOpen$, $disconnected$, $reconnecting$, or $closed$
    connection has no valid handles.  This generalises DES-029 and
    is the safety half of the biff-tww fix: stale handles cached
    during a wedge are, in the model, no handles at all.

  \item \textbf{Wedge liveness (biff-tww)}: no reachable state is
    stuck in $halfOpen$.  From $halfOpen$ the recovery sequence
    $WedgeDetected \to Reconnect$ (via $reconnecting$) is always
    enabled, so the connection cannot deadlock in the wedge.  The
    proof is by deadlock analysis: removing $WedgeDetected$ (the
    pre-fix code, whose keepalive ping never fired in time) makes
    $halfOpen$ a reachable terminal state.  probcli surfaces it as
    a deadlock along the trace
    $EnsureConnected \to Request(reqTimeout)$, which \emph{is}
    biff-tww.  With $WedgeDetected$ present the state space is
    deadlock-free.

  \item \textbf{Snapshot survival}: $SnapshotDone$ transitions to
    $streaming$, never to $idle$ or $stopped$.  The watcher
    remains alive.

  \item \textbf{Timeout resilience}: $HandleTimeout$ keeps the
    watcher in $streaming$ phase.  No restart, no phase change.

  \item \textbf{Error restart}: Only $HandleError$ increments
    $restartCount$.  Snapshot-done and timeout do not.

  \item \textbf{Connection prerequisite}: The watcher cannot be
    active ($phase \neq idle$) when the connection is not
    $connected$; every session and watcher operation is enabled
    only while $connState = connected$.

  \item \textbf{Session uniqueness}: The partial injection
    ($\pinj$) on sessions guarantees each session key maps to
    exactly one value.
\end{enumerate}

\end{document}
